\documentclass[10pt,a4paper]{amsart}
\usepackage[margin=19mm]{geometry}
\usepackage{amsmath,amssymb,mathtools}
\usepackage[scr=rsfs,cal=euler]{mathalfa}
\usepackage{xfrac}
\usepackage[unicode,hidelinks,bookmarks=false]{hyperref}
\newcommand{\ie}{i.e.\ }
\newcommand{\cf}{cf.\ }
\newcommand{\resp}{respectively\ }
\newcommand{\Subsection}{\subsection}
\providecommand{\nobreakdash}{\nobreak\hbox{-}}
\setlength{\emergencystretch}{2em}
\multlinegap0pt
\usepackage{bm}
\usepackage{bbm}
\usepackage{dsfont}
\usepackage{xspace}
\usepackage{cancel}
\usepackage{mathtools}
\usepackage{amsthm}
\makeatletter
\@ifundefined{lemma}{\newtheorem{lemma}{Lemma}}{}
\makeatother
\title[The volume comparison of symmetric spaces of non-compact typ]{The volume comparison of symmetric spaces of non-compact type of rank 1}
\author{Jiaqi Chen, Yufei Shan, Yinghui Ye}
\date{Source: arXiv:2508.08674v2 (CC BY 4.0)}
\begin{document}
\maketitle

\noindent\emph{Source and licence.} The volume comparison of symmetric spaces of non-compact type of rank 1. Version arXiv:2508.08674v2 (CC BY 4.0), distributed under the Creative Commons Attribution 4.0 International licence, as recorded in the arXiv licence metadata for this exact version and shown on the abstract page https://arxiv.org/abs/2508.08674v2. Full rights notice: licenses/arxiv-2508.08674-CC-BY-4.0.txt.\par\smallskip

\noindent\textit{Editorial scope.} Counted unit: the Lemma whose source label is \texttt{spherical frame} in the article (source0.tex lines 268--289, source label \texttt{spherical frame}), delivered together with the article's own complete original proof of it (source0.tex lines 290--365, one \texttt{proof} environment of 76 lines). No mathematics is written, repaired or re-derived: statement and proof are the article's own text line for line. Required context retained uncounted: relation between lie algebra and kill field (lines 161--163); metric = Killing form (lines 183--185); sectional curvature (lines 187--192); inner product on lie algebra (lines 212--214); sinh existence (lines 236--241); basic frame relation 1 (lines 263--265). Every definition, display and label of those lines is retained unchanged. Omitted from this package and not redistributed: 1--160, 164--182, 186--186, 193--211, 215--235, 242--262, 266--267, 366--1612. Nothing inside a retained excerpt is omitted or altered: every retained body is delivered exactly as the article's lines are written, so no omission receipt of any kind accompanies this package. Citations: the retained excerpts carry no citation command at all, so no bibliography entry of the article is transcribed and no reference is redistributed. Reference adaptation: none was needed and none was made; every reference inside the delivered unit resolves inside it, so no unresolved reference or citation is produced. Printed slips kept literal: no printed word, symbol, hypothesis or number of the counted statement or of its proof was altered. Layout and class: the article's own class is \texttt{amsart} with packages including fontenc, inputenc, fullpage, appendix, verbatim, amsmath, amssymb, float, color, bbding, wrapfig, pst-node, biblatex, lineno; the wrapper is the lane's pinned \texttt{amsart} editorial wrapper at 10pt on A4, which restates the theorem-like environments the retained text needs and adds the article's own macro definitions that the retained text needs (none), with extra wrapper packages (bm, bbm, dsfont, xspace, cancel, mathtools). The delivered numbering is the unit's own. Assets: no retained span contains a figure environment, a graphics-inclusion command, a PDF-page inclusion, a table or a TikZ picture; no image file is copied into this package, the package declares no asset dependency and no image file is redistributed with it. Licence: version arXiv:2508.08674v2 of this article is distributed under the Creative Commons Attribution 4.0 International licence, as recorded in the arXiv licence metadata for this exact version and shown on the abstract page https://arxiv.org/abs/2508.08674v2. A full rights notice is delivered at licenses/arxiv-2508.08674-CC-BY-4.0.txt. Characters: every retained excerpt is pure ASCII, so the delivered engine, XeLaTeX with its default Latin Modern text font, reports no missing character.
\begin{align}\label{relation between lie algebra and kill field}
    \sigma([x, y]) = -[\sigma(x), \sigma(y)], 
\end{align}

\begin{align}\label{metric = Killing form}
g_0\bigl(\sigma(x), \sigma(y)\bigr)\big|_{p_0} = \lambda \langle x, y\rangle,
\end{align}

\begin{equation}\label{sectional curvature}
    \begin{split}
        (\nabla_{\hat{X}} Y)|_{p_0} =& [\hat{X}, Y]|_{p_0},\\
        R(\hat{X}, \hat{Y}) \hat{Z}|_{p_0} =& -[[\hat{X}, \hat{Y}], \hat{Z}]|_{p_0} = -\sigma([[x, y], z])|_{p_0},
    \end{split}
\end{equation}

\begin{align}\label{inner product on lie algebra}
	(\cdot, \cdot): = -\lambda \langle \cdot, d\xi(\cdot) \rangle,
\end{align}

\begin{equation}\label{sinh existence}
    \begin{split}
        \exp(ad(v))(p_i) =& \mathrm{sh}(a_i(v))k_i + \mathrm{ch}(a_i(v)) p_i,\\
        \exp(ad(v))(k_i) =& \mathrm{sh}(a_i(v))p_i + \mathrm{ch}(a_i(v)) k_i,
    \end{split}
\end{equation}

\begin{equation}\label{basic frame relation 1}
    [p_i, k_i] = \frac{1}{2}[x_i - y_i, x_i + y_i] = [x_i,y_i] = -\sum_{k = 1}^{r}a_i(v_k)v_k,
\end{equation}

\begin{lemma}\label{spherical frame}
    Let $(M^n, g_0)$ be a simply connected irreducible symmetric space of non-compact type of rank $r$. Then, under the above notation, for any $v\in \mathfrak{a}$, we have that 
    \begin{equation}\label{orthonormal basis on radius}
        \begin{split}
            &g_0(\sigma(v_i), \sigma(k_j))|_{\exp(v t)(p_0)} = 0,\\
            &g_0(\sigma(v_{i_1}),\sigma(v_{i_2}))|_{\exp(vt)(p_0)} = \delta_{i_1i_2},\\
            &g_0(\sigma(k_{j_1}), \sigma(k_{j_2}))|_{\exp(vt)(p_0)} = \mathrm{sh}^2(a_{j_1}(v)t)\delta_{j_1j_2},
        \end{split}
    \end{equation}
    for any $1\leq i, i_1, i_2 \leq r$ and $1 \leq j, j_1, j_2 \leq n - r$. Moreover, we have that 
    \begin{itemize}
        \item for $1 \leq i \leq r$,
        \begin{equation}\label{covariant derivative on radius 1}
            \nabla_{\sigma(v_i)}\sigma(v_i)|_{\exp(vt)(p_0)} = 0,
        \end{equation}
        \item for $1 \leq j \leq n - r$,
        \begin{equation}\label{covariant derivative on radius 2}
            \nabla_{\sigma(k_j)}\sigma(k_j)|_{\exp(vt)(p_0)}\\ 
            = \mathrm{sh}(-a_j(v)t)\cdot \mathrm{ch}(-a_j(v)t)\cdot \sum_{l = 1}^{r}a_j(v_l)\sigma(v_l)|_{\exp(vt)(p_0)}. 
        \end{equation}
    \end{itemize}
\end{lemma}

\begin{proof}
    For \eqref{orthonormal basis on radius}, since $\exp(-v t)$ is in the isometry group of $M$ for any $v\in \mathfrak{a}$, it suffices to show that 
    \begin{align*}
        &g_0(\exp(-vt)_*(\sigma(v_i)|_{\exp(vt)(p_0)}),\exp(-vt)_*(\sigma(k_j)|_{\exp(vt)(p_0)})) = 0,\\
        &g_0(\exp(-vt)_*(\sigma(v_{i_1})|_{\exp(vt)(p_0)}),\exp(-vt)_*(\sigma(v_{i_2})|_{\exp(vt)(p_0)})) = \delta_{i_1i_2},\\
        &g_0(\exp(-vt)_*(\sigma(k_{j_1})|_{\exp(vt)(p_0)}),\exp(-vt)_*(\sigma(k_{j_2})|_{\exp(vt)(p_0)})) = \mathrm{sh}^2(a_{j_1}(v)t)\delta_{j_1j_2}, 
    \end{align*}
    for $1 \leq i, i_1, i_2 \leq r$ and $1 \leq j, j_1, j_2 \leq n - r$. In fact, for $1 \leq i \leq r$, 
    \begin{align*}
        \exp(-vt)_*(\sigma(v_i)|_{\exp(vt)(p_0)}) 
        =& \exp(-vt)_*(\frac{d}{ds} \bigg|_{s = 0} \exp(v_is)\cdot\exp(vt)(p_0))\\
        =& \frac{d}{ds} \bigg|_{s = 0} \exp(-vt)\cdot\exp(v_is)\cdot \exp(vt)(p_0)\\
        =& \frac{d}{ds} \bigg|_{s = 0} \exp([\exp(ad(-vt))v_i]s)\\
        =& \sigma(\exp(ad(-vt))v_i)|_{p_0}\\
        =& \sigma(v_i)|_{p_0},
    \end{align*}
    and for $1\leq j \leq n - r$, by \eqref{sinh existence}, we have that 
    \begin{align*}
        \exp(-vt)_*(\sigma(k_j)|_{\exp(vt)(p_0)}) 
        =& \exp(-vt)_*(\frac{d}{ds}  \bigg|_{s = 0}  \exp(k_js) \cdot \exp(vt)(p_0))\\
        =& \frac{d}{ds} \bigg|_{s = 0} \exp(-vt)\cdot\exp(k_js)\cdot\exp(vt)(p_0)\\
        =& \frac{d}{ds} \bigg|_{s = 0} \exp([\exp(ad(-vt))(k_i)]s)(p_0)\\
        =& \frac{d}{ds} \bigg|_{s = 0} \exp([\mathrm{sh}(-a_i(v)t)p_i + \mathrm{ch}(-a_i(v)t)k_i]s)(p_0)\\
        =& \sigma(\mathrm{sh}(-a_i(v)t)p_i + \mathrm{ch}(-a_i(v)t)k_i)|_{p_0}\\
        =& \mathrm{sh}(-a_i(v)t)\sigma(p_i)|_{p_0}. 
    \end{align*}
    Then, together with \eqref{metric = Killing form} and \eqref{inner product on lie algebra}, \eqref{orthonormal basis on radius}) follows. 

    For \eqref{covariant derivative on radius 1} and \eqref{covariant derivative on radius 2}, let $1 \leq j \leq n - r$. The integral curve $\gamma_j(s)$ of the Killing field $\sigma(k_j)$ passing through $\exp(vt)(p_0)$ is $\gamma_j(s) := \exp(k_js)\cdot \exp(vt)(p_0)$. Then, $\sigma(k_j)|_{\gamma_j(s)} = \dot\gamma_j(s)$. Thus, 
    \begin{align*}
        \nabla_{\sigma(k_j)}\sigma(k_j)|_{\exp(vt)(p_0)} = \nabla_{\dot\gamma_j(s)}\dot\gamma_j(s)|_{s = 0}.
    \end{align*}
    Then, since $\exp(-vt)$ is an isometry with its inverse $\exp(vt)$, we have that 
    \begin{align*}
        \nabla_{\dot\gamma_j(s)}\dot\gamma_j(s)|_{s = 0} = \exp(vt)_* (\nabla_{\exp(-vt)_*(\dot\gamma_j(s))}\exp(-vt)_*(\dot\gamma_j(s))|_{p_0}).
    \end{align*}
    On the other hand, by \eqref{sinh existence},
    \begin{align*}
        \exp(-vt)_*(\dot\gamma_j(s)) 
        =& \frac{d}{ds}\exp(-vt)(\gamma_j(s))\\
        =& \frac{d}{ds}\exp(-vt)\exp(k_j s)\exp(vt)(p_0)\\
        =& \frac{d}{ds} \exp([\exp(ad(-vt))(k_j)]s)(p_0)\\
        =& \frac{d}{ds} \exp([\mathrm{sh}(-a_j(v)t)p_j + \mathrm{ch}(-a_j(v)t)k_j]s)(p_0)\\
        =& \sigma(\mathrm{sh}(-a_j(v)t)p_j + \mathrm{ch}(-a_j(v)t)k_j)|_{exp(-vt)(\gamma_j(s))},
    \end{align*}
    and at $p_0 \in M$, 
    \begin{align*}
        \exp(-vt)_*(\dot\gamma_j(s))|_{s = 0} = \sigma(\mathrm{sh}(-a_j(v)t)p_j + \mathrm{ch}(-a_j(v)t)k_j)|_{p_0} = \mathrm{sh}(-a_j(v)t)\cdot\sigma(p_j)|_{p_0}.
    \end{align*}
    Therefore, by \eqref{relation between lie algebra and kill field}, \eqref{sectional curvature} and \eqref{basic frame relation 1}, we have  
    \begin{align*}
        &\nabla_{\exp(-vt)_*(\dot\gamma_j(s))}\exp(-vt)_*(\dot\gamma_j(s))|_{p_0}\\
        =& \mathrm{sh}(-a_j(v)t)\cdot\nabla_{\sigma(p_j)}\sigma(\mathrm{sh}(-a_j(v)t)p_j + \mathrm{ch}(-a_j(v)t)k_j)|_{p_0}\\
        =& \mathrm{sh}(-a_j(v)t)\cdot\Big[\sigma(p_j), \mathrm{sh}(-a_j(v)t)\sigma(p_j) + \mathrm{ch}(-a_j(v)t)\sigma(k_j)\Big]\Big|_{p_0}\\
        =& -\mathrm{sh}(-a_j(v)t)\cdot \mathrm{ch}(-a_j(v)t)\cdot\sigma([p_j, k_j])|_{p_0}\\
        =& \mathrm{sh}(-a_j(v)t)\cdot \mathrm{ch}(-a_j(v)t)\cdot \sum_{k = 0}^{r}a_j(v_k)\sigma(v_k)|_{p_0}. 
    \end{align*}
    Thereby, 
    \begin{align*}
        &\exp(vt)_*(\nabla_{\exp(-vt)_*(\dot\gamma_j(s))}\exp(-vt)_*(\dot\gamma_j(s))|_{p_0})\\
        =& \mathrm{sh}(-a_j(v)t)\cdot \mathrm{ch}(-a_j(v)t)\cdot \sum_{k = 0}^{r}a_j(v_k)\sigma(v_k)|_{\exp(vt)(p_0)}.
    \end{align*}
    The last step is followed by 
    \begin{align*}
        \exp(vt)_*(\sigma(v_k)|_{p_0}) 
        =& \frac{d}{ds} \bigg|_{s = 0} \exp(vt)\cdot \exp(v_ks)(p_0)\\
        =& \frac{d}{ds} \bigg|_{s = 0} \exp(vt)\cdot \exp(v_ks)\cdot \exp(-vt)\cdot \exp(vt)(p_0)\\
        =& \frac{d}{ds} \bigg|_{s = 0} \exp([\exp(ad(vt))v_k]s)(\exp(vt)(p_0))\\
        =& \frac{d}{ds} \bigg|_{s = 0} \exp(v_ks)(\exp(vt)(p_0))\\
        =& \sigma(v_k)|_{\exp(vt)(p_0)}.
    \end{align*}
    By the similar steps, we can get that for $1 \leq i \leq r$, 
    \begin{align*}
        \nabla_{\sigma(v_i)}\sigma(v_i)|_{\exp(vt)(p_0)} = 0.
    \end{align*}
\end{proof}

\end{document}
